% =====================================================================
% Paper 05 — Pillar 5: Chirality from shell-genericity, opposite-valley
%   pairing, mass protection, and the reduced Weyl/Dirac witness theorem
%   on the BCC condensate. (Rev 3, 2026-05-03)
% Status: [DRAFT] [NEEDS_UPDATE] — Wave 3 / Pillar 5
% AUDIT-FLAG: AUDIT-2026-05-02-Wave7-Aux-Epoch-Overclaim
%   Rev 1 (2026-05-02): downgraded to "mechanism note" per first-pass
%     audit, but body kept Index(D)=wind(m) framework + T7 PROVED
%     wording, both flagged as defective by the audit.
%   Rev 2 (2026-05-02): no real fix; audit identified that Rev 1
%     was effectively a cover-page change only.
%   Rev 3 (2026-05-03): COMPLETE REWRITE based on the actual canonical
%     Math10-14 archive, replacing the incorrect Index(D)=wind(m)
%     framework with the canonical Math10-14 + Math171-AddA mechanism:
%     (i) Math10 shell-genericity (Generic Nonvanishing of Shell-
%         Protected Dirac Couplings; Single-Channel + Two-Channel
%         Genericity theorems; Dirac Genericity Principle);
%     (ii) Math11 minimal-model uniqueness (five constraints forcing
%         the complex triplet);
%     (iii) Math12 mass protection (classification of constant Dirac
%         masses on the doubled block; only specific terms open a
%         genuine Dirac gap);
%     (iv) Math13 local Dirac block reduction
%         (H_D = C^2_valley x C^2_int);
%     (v) Math14 reduced Weyl/Dirac witness theorems (three theorems:
%         linearised quartic no-go, exact shell cancellation, locked
%         cross-shell quartic vanishing — closing the reduced witness
%         in Class II seed space);
%     (vi) Math171-AddA HRR formula ind(D_E^c) = r - μ (per Top-impact
%         paper TI-1) for the family-count specialisation r=16 -> μ=0.
%   Note: the previously incorrect Index(D)=wind(m) on a 3D Brillouin
%   torus (without Callias/domain-wall structure) is REMOVED; the
%   correct framework is the one above. A Callias-type rigorous
%   theorem-level rewrite remains an open programme item per
%   Q-2026-05-02-Pillar5-Mechanism-Rewrite.
% Compile: pdflatex Paper-05.tex (×2 for refs)
% Canonical archive: Math10, Math11, Math12, Math13, Math14, Math171-AddA,
%   Math233 (spinor pair-kernel programme), Math314-AddD (audit)
% =====================================================================
% --- TECT canonical preamble (see _shared/tect-preamble.tex) ----------
% =====================================================================
% TECT canonical LaTeX preamble (revtex4-2 / single-column / theorem-safe)
% =====================================================================
% Maintainer: Jusang Lee (jtkor@outlook.com)
% Binding from: 2026-05-06
% Purpose: a single source-of-truth preamble for every TECT paper
% (Paper-00 .. Paper-10 + Auxiliary notes). Designed to (a) eliminate
% font-corruption on pdflatex (T1 fontenc + lmodern), (b) make
% theorem-heavy mathematical-physics content render cleanly in a
% single-column layout (PRD class option, NOT PRL two-column), (c)
% provide consistent theorem numbering across all TECT papers via
% amsthm with a shared counter set, (d) make hyperref + cleveref
% cross-references resolve cleanly with the standard two-pass
% pdflatex compile.
%
% Usage in each Paper-NN.tex:
%   % =====================================================================
% TECT canonical LaTeX preamble (revtex4-2 / single-column / theorem-safe)
% =====================================================================
% Maintainer: Jusang Lee (jtkor@outlook.com)
% Binding from: 2026-05-06
% Purpose: a single source-of-truth preamble for every TECT paper
% (Paper-00 .. Paper-10 + Auxiliary notes). Designed to (a) eliminate
% font-corruption on pdflatex (T1 fontenc + lmodern), (b) make
% theorem-heavy mathematical-physics content render cleanly in a
% single-column layout (PRD class option, NOT PRL two-column), (c)
% provide consistent theorem numbering across all TECT papers via
% amsthm with a shared counter set, (d) make hyperref + cleveref
% cross-references resolve cleanly with the standard two-pass
% pdflatex compile.
%
% Usage in each Paper-NN.tex:
%   % =====================================================================
% TECT canonical LaTeX preamble (revtex4-2 / single-column / theorem-safe)
% =====================================================================
% Maintainer: Jusang Lee (jtkor@outlook.com)
% Binding from: 2026-05-06
% Purpose: a single source-of-truth preamble for every TECT paper
% (Paper-00 .. Paper-10 + Auxiliary notes). Designed to (a) eliminate
% font-corruption on pdflatex (T1 fontenc + lmodern), (b) make
% theorem-heavy mathematical-physics content render cleanly in a
% single-column layout (PRD class option, NOT PRL two-column), (c)
% provide consistent theorem numbering across all TECT papers via
% amsthm with a shared counter set, (d) make hyperref + cleveref
% cross-references resolve cleanly with the standard two-pass
% pdflatex compile.
%
% Usage in each Paper-NN.tex:
%   \input{../_shared/tect-preamble.tex}
%   \begin{document}
%   ...
%   \end{document}
%
% Build (every paper): `pdflatex Paper-NN && pdflatex Paper-NN`
% (the second pass resolves \ref{} / \cref{} forward references).
% A bash/PowerShell wrapper that automates this is at
% Docs/papers/papers/_shared/build-paper.sh / build-paper.ps1.
% =====================================================================

% ---- Document class --------------------------------------------------
% PRD class option of revtex4-2 with [reprint,onecolumn] gives the same
% APS heading / by-line / abstract style as PRL but in a wider single
% column that handles long display equations and theorem environments
% without column-break artefacts. nofootinbib keeps citations out of
% footnotes. The 11pt option restores standard font metrics; with T1
% fontenc + lmodern this gives non-bitmap glyphs.
\documentclass[%
  aps,prd,reprint,onecolumn,nofootinbib,11pt,%
  amsmath,amssymb,floatfix,longbibliography%
]{revtex4-2}

% ---- Encoding & fonts ------------------------------------------------
\usepackage[T1]{fontenc}        % proper 8-bit font encoding (no font corruption)
\usepackage[utf8]{inputenc}     % allow UTF-8 source
\usepackage{lmodern}            % scalable Latin Modern (replaces bitmap CM)
\usepackage{textcomp}           % extra text symbols
\usepackage{microtype}          % subtle kerning improvements

% ---- UTF-8 → LaTeX character mappings --------------------------------
% Maps the Unicode characters that occur in TECT body text (legacy from
% the chat-content auto-archival rule, CLAUDE.md §4) to LaTeX-safe
% commands. Without these declarations, T1+inputenc rejects the chars
% with "Unicode character X not set up for use with LaTeX".
% Adding declarations centrally (here) is preferable to per-file edits.
\DeclareUnicodeCharacter{00A7}{\S}              % §  (section sign)
\DeclareUnicodeCharacter{00B2}{\ensuremath{^{2}}}        % ²
\DeclareUnicodeCharacter{00B3}{\ensuremath{^{3}}}        % ³
\DeclareUnicodeCharacter{00B5}{\ensuremath{\mu}}         % µ (micro)
\DeclareUnicodeCharacter{00B7}{\ensuremath{\cdot}}       % · (middle dot)
\DeclareUnicodeCharacter{00D7}{\ensuremath{\times}}      % ×
\DeclareUnicodeCharacter{00E4}{\"a}             % ä
\DeclareUnicodeCharacter{00E9}{\'e}             % é
\DeclareUnicodeCharacter{00F6}{\"o}             % ö
\DeclareUnicodeCharacter{00FC}{\"u}             % ü
\DeclareUnicodeCharacter{010C}{\v{C}}           % Č
\DeclareUnicodeCharacter{010D}{\v{c}}           % č
\DeclareUnicodeCharacter{0394}{\ensuremath{\Delta}}      % Δ
\DeclareUnicodeCharacter{03A3}{\ensuremath{\Sigma}}      % Σ
\DeclareUnicodeCharacter{03A9}{\ensuremath{\Omega}}      % Ω
\DeclareUnicodeCharacter{03B1}{\ensuremath{\alpha}}      % α
\DeclareUnicodeCharacter{03B2}{\ensuremath{\beta}}       % β
\DeclareUnicodeCharacter{03B3}{\ensuremath{\gamma}}      % γ
\DeclareUnicodeCharacter{03B4}{\ensuremath{\delta}}      % δ
\DeclareUnicodeCharacter{03B5}{\ensuremath{\varepsilon}} % ε
\DeclareUnicodeCharacter{03BA}{\ensuremath{\kappa}}      % κ
\DeclareUnicodeCharacter{03BB}{\ensuremath{\lambda}}     % λ
\DeclareUnicodeCharacter{03BC}{\ensuremath{\mu}}         % μ
\DeclareUnicodeCharacter{03BD}{\ensuremath{\nu}}         % ν
\DeclareUnicodeCharacter{03C0}{\ensuremath{\pi}}         % π
\DeclareUnicodeCharacter{03C1}{\ensuremath{\rho}}        % ρ
\DeclareUnicodeCharacter{03C3}{\ensuremath{\sigma}}      % σ
\DeclareUnicodeCharacter{03C4}{\ensuremath{\tau}}        % τ
\DeclareUnicodeCharacter{03C6}{\ensuremath{\varphi}}     % φ
\DeclareUnicodeCharacter{03C7}{\ensuremath{\chi}}        % χ
\DeclareUnicodeCharacter{03C8}{\ensuremath{\psi}}        % ψ
\DeclareUnicodeCharacter{03C9}{\ensuremath{\omega}}      % ω
\DeclareUnicodeCharacter{2013}{--}              % – (en dash)
\DeclareUnicodeCharacter{2014}{---}             % — (em dash)
\DeclareUnicodeCharacter{2018}{`}               % ‘
\DeclareUnicodeCharacter{2019}{'}               % ’
\DeclareUnicodeCharacter{201C}{``}              % “
\DeclareUnicodeCharacter{201D}{''}              % ”
\DeclareUnicodeCharacter{2070}{\ensuremath{^{0}}}        % ⁰
\DeclareUnicodeCharacter{2071}{\ensuremath{^{i}}}        % ⁱ
\DeclareUnicodeCharacter{2122}{\textsuperscript{TM}}     % ™
\DeclareUnicodeCharacter{2126}{\ensuremath{\Omega}}      % Ω (ohm sign)
\DeclareUnicodeCharacter{210F}{\ensuremath{\hbar}}       % ℏ
\DeclareUnicodeCharacter{2115}{\ensuremath{\mathbb{N}}}  % ℕ
\DeclareUnicodeCharacter{2102}{\ensuremath{\mathbb{C}}}  % ℂ
\DeclareUnicodeCharacter{211D}{\ensuremath{\mathbb{R}}}  % ℝ
\DeclareUnicodeCharacter{2124}{\ensuremath{\mathbb{Z}}}  % ℤ
\DeclareUnicodeCharacter{211A}{\ensuremath{\mathbb{Q}}}  % ℚ
\DeclareUnicodeCharacter{2192}{\ensuremath{\to}}         % →
\DeclareUnicodeCharacter{21D2}{\ensuremath{\Rightarrow}} % ⇒
\DeclareUnicodeCharacter{2200}{\ensuremath{\forall}}     % ∀
\DeclareUnicodeCharacter{2203}{\ensuremath{\exists}}     % ∃
\DeclareUnicodeCharacter{2208}{\ensuremath{\in}}         % ∈
\DeclareUnicodeCharacter{2209}{\ensuremath{\notin}}      % ∉
\DeclareUnicodeCharacter{2227}{\ensuremath{\wedge}}      % ∧
\DeclareUnicodeCharacter{2228}{\ensuremath{\vee}}        % ∨
\DeclareUnicodeCharacter{2229}{\ensuremath{\cap}}        % ∩
\DeclareUnicodeCharacter{222A}{\ensuremath{\cup}}        % ∪
\DeclareUnicodeCharacter{2245}{\ensuremath{\cong}}       % ≅
\DeclareUnicodeCharacter{2248}{\ensuremath{\approx}}     % ≈
\DeclareUnicodeCharacter{2260}{\ensuremath{\neq}}        % ≠
\DeclareUnicodeCharacter{2264}{\ensuremath{\leq}}        % ≤
\DeclareUnicodeCharacter{2265}{\ensuremath{\geq}}        % ≥
\DeclareUnicodeCharacter{22C5}{\ensuremath{\cdot}}       % ⋅
\DeclareUnicodeCharacter{2295}{\ensuremath{\oplus}}      % ⊕
\DeclareUnicodeCharacter{2297}{\ensuremath{\otimes}}     % ⊗

% ---- Mathematics & symbols ------------------------------------------
\usepackage{amsmath,amssymb,bm,mathrsfs,mathtools}
\usepackage{amsthm}             % theorem environments (load BEFORE hyperref)

% ---- Theorem environments (shared counter set, TECT canonical) -------
% All theorems / lemmas / propositions / corollaries / remarks share a
% single counter so cross-references read consistently across a paper.
% Definitions get a separate counter (definitions are numbered in their
% own sequence by editorial convention).
\newtheorem{theorem}{Theorem}
\newtheorem{lemma}[theorem]{Lemma}
\newtheorem{proposition}[theorem]{Proposition}
\newtheorem{corollary}[theorem]{Corollary}

\theoremstyle{remark}
\newtheorem{remark}[theorem]{Remark}
\newtheorem{example}[theorem]{Example}

\theoremstyle{definition}
\newtheorem{definition}[theorem]{Definition}
\newtheorem{conjecture}[theorem]{Conjecture}
\newtheorem{hypothesis}[theorem]{Hypothesis}

% ---- Tables / floats / graphics --------------------------------------
\usepackage{booktabs}           % \toprule / \midrule / \bottomrule
\usepackage{array}
\usepackage{graphicx}
\usepackage{enumitem}           % \begin{itemize}[leftmargin=...] options
\usepackage{hyphenat}           % \hyp{} for breakable hyphens in \texttt
\usepackage{seqsplit}           % \seqsplit{} for long unbreakable strings
\sloppy                         % allow stretchier inter-word spacing to
                                % reduce overfull \hbox warnings on long URLs
                                % and long \texttt tokens (paragraph-level).
\emergencystretch=3em           % last-resort hbox stretch budget

% ---- Cross-references (hyperref last among the standard packages,
%      cleveref AFTER hyperref) ----------------------------------------
\usepackage[%
  colorlinks=true,%
  linkcolor=blue,%
  citecolor=blue,%
  urlcolor=blue,%
  breaklinks=true,%
  bookmarks=true,%
  bookmarksnumbered=true%
]{hyperref}
\usepackage[capitalise,nameinlink]{cleveref}

% ---- TECT-canonical author block macros -----------------------------
% (defined here so every paper has the same author / affiliation style)
\newcommand{\tectauthor}{%
  \author{Jusang Lee}%
  \email{jtkor@outlook.com}%
  \affiliation{Independent researcher, Republic of Korea}%
}

% ---- TECT-canonical archive footer ----------------------------------
\newcommand{\tectarchivefooter}{%
  \par\medskip
  \noindent\small\textit{Canonical archive}:
  \url{https://github.com/TECT-OpenPhysics/TECT}.\par%
}

% ---- TECT TODO / AUDIT-FLAG / HONEST-SCOPE inline marks --------------
% Used in drafts to highlight pending audit items without disturbing
% the body text style. Suppress via \tectsuppressmarks (define empty).
\providecommand{\tectauditflag}[1]{\textcolor{red}{\textbf{[AUDIT-FLAG: #1]}}}
\providecommand{\tecttodo}[1]{\textcolor{orange}{\textbf{[TODO: #1]}}}
\providecommand{\tecthonestscope}[1]{\textit{Honest scope: #1.}}

% =====================================================================
% End of TECT canonical preamble
% =====================================================================

%   \begin{document}
%   ...
%   \end{document}
%
% Build (every paper): `pdflatex Paper-NN && pdflatex Paper-NN`
% (the second pass resolves \ref{} / \cref{} forward references).
% A bash/PowerShell wrapper that automates this is at
% Docs/papers/papers/_shared/build-paper.sh / build-paper.ps1.
% =====================================================================

% ---- Document class --------------------------------------------------
% PRD class option of revtex4-2 with [reprint,onecolumn] gives the same
% APS heading / by-line / abstract style as PRL but in a wider single
% column that handles long display equations and theorem environments
% without column-break artefacts. nofootinbib keeps citations out of
% footnotes. The 11pt option restores standard font metrics; with T1
% fontenc + lmodern this gives non-bitmap glyphs.
\documentclass[%
  aps,prd,reprint,onecolumn,nofootinbib,11pt,%
  amsmath,amssymb,floatfix,longbibliography%
]{revtex4-2}

% ---- Encoding & fonts ------------------------------------------------
\usepackage[T1]{fontenc}        % proper 8-bit font encoding (no font corruption)
\usepackage[utf8]{inputenc}     % allow UTF-8 source
\usepackage{lmodern}            % scalable Latin Modern (replaces bitmap CM)
\usepackage{textcomp}           % extra text symbols
\usepackage{microtype}          % subtle kerning improvements

% ---- UTF-8 → LaTeX character mappings --------------------------------
% Maps the Unicode characters that occur in TECT body text (legacy from
% the chat-content auto-archival rule, CLAUDE.md §4) to LaTeX-safe
% commands. Without these declarations, T1+inputenc rejects the chars
% with "Unicode character X not set up for use with LaTeX".
% Adding declarations centrally (here) is preferable to per-file edits.
\DeclareUnicodeCharacter{00A7}{\S}              % §  (section sign)
\DeclareUnicodeCharacter{00B2}{\ensuremath{^{2}}}        % ²
\DeclareUnicodeCharacter{00B3}{\ensuremath{^{3}}}        % ³
\DeclareUnicodeCharacter{00B5}{\ensuremath{\mu}}         % µ (micro)
\DeclareUnicodeCharacter{00B7}{\ensuremath{\cdot}}       % · (middle dot)
\DeclareUnicodeCharacter{00D7}{\ensuremath{\times}}      % ×
\DeclareUnicodeCharacter{00E4}{\"a}             % ä
\DeclareUnicodeCharacter{00E9}{\'e}             % é
\DeclareUnicodeCharacter{00F6}{\"o}             % ö
\DeclareUnicodeCharacter{00FC}{\"u}             % ü
\DeclareUnicodeCharacter{010C}{\v{C}}           % Č
\DeclareUnicodeCharacter{010D}{\v{c}}           % č
\DeclareUnicodeCharacter{0394}{\ensuremath{\Delta}}      % Δ
\DeclareUnicodeCharacter{03A3}{\ensuremath{\Sigma}}      % Σ
\DeclareUnicodeCharacter{03A9}{\ensuremath{\Omega}}      % Ω
\DeclareUnicodeCharacter{03B1}{\ensuremath{\alpha}}      % α
\DeclareUnicodeCharacter{03B2}{\ensuremath{\beta}}       % β
\DeclareUnicodeCharacter{03B3}{\ensuremath{\gamma}}      % γ
\DeclareUnicodeCharacter{03B4}{\ensuremath{\delta}}      % δ
\DeclareUnicodeCharacter{03B5}{\ensuremath{\varepsilon}} % ε
\DeclareUnicodeCharacter{03BA}{\ensuremath{\kappa}}      % κ
\DeclareUnicodeCharacter{03BB}{\ensuremath{\lambda}}     % λ
\DeclareUnicodeCharacter{03BC}{\ensuremath{\mu}}         % μ
\DeclareUnicodeCharacter{03BD}{\ensuremath{\nu}}         % ν
\DeclareUnicodeCharacter{03C0}{\ensuremath{\pi}}         % π
\DeclareUnicodeCharacter{03C1}{\ensuremath{\rho}}        % ρ
\DeclareUnicodeCharacter{03C3}{\ensuremath{\sigma}}      % σ
\DeclareUnicodeCharacter{03C4}{\ensuremath{\tau}}        % τ
\DeclareUnicodeCharacter{03C6}{\ensuremath{\varphi}}     % φ
\DeclareUnicodeCharacter{03C7}{\ensuremath{\chi}}        % χ
\DeclareUnicodeCharacter{03C8}{\ensuremath{\psi}}        % ψ
\DeclareUnicodeCharacter{03C9}{\ensuremath{\omega}}      % ω
\DeclareUnicodeCharacter{2013}{--}              % – (en dash)
\DeclareUnicodeCharacter{2014}{---}             % — (em dash)
\DeclareUnicodeCharacter{2018}{`}               % ‘
\DeclareUnicodeCharacter{2019}{'}               % ’
\DeclareUnicodeCharacter{201C}{``}              % “
\DeclareUnicodeCharacter{201D}{''}              % ”
\DeclareUnicodeCharacter{2070}{\ensuremath{^{0}}}        % ⁰
\DeclareUnicodeCharacter{2071}{\ensuremath{^{i}}}        % ⁱ
\DeclareUnicodeCharacter{2122}{\textsuperscript{TM}}     % ™
\DeclareUnicodeCharacter{2126}{\ensuremath{\Omega}}      % Ω (ohm sign)
\DeclareUnicodeCharacter{210F}{\ensuremath{\hbar}}       % ℏ
\DeclareUnicodeCharacter{2115}{\ensuremath{\mathbb{N}}}  % ℕ
\DeclareUnicodeCharacter{2102}{\ensuremath{\mathbb{C}}}  % ℂ
\DeclareUnicodeCharacter{211D}{\ensuremath{\mathbb{R}}}  % ℝ
\DeclareUnicodeCharacter{2124}{\ensuremath{\mathbb{Z}}}  % ℤ
\DeclareUnicodeCharacter{211A}{\ensuremath{\mathbb{Q}}}  % ℚ
\DeclareUnicodeCharacter{2192}{\ensuremath{\to}}         % →
\DeclareUnicodeCharacter{21D2}{\ensuremath{\Rightarrow}} % ⇒
\DeclareUnicodeCharacter{2200}{\ensuremath{\forall}}     % ∀
\DeclareUnicodeCharacter{2203}{\ensuremath{\exists}}     % ∃
\DeclareUnicodeCharacter{2208}{\ensuremath{\in}}         % ∈
\DeclareUnicodeCharacter{2209}{\ensuremath{\notin}}      % ∉
\DeclareUnicodeCharacter{2227}{\ensuremath{\wedge}}      % ∧
\DeclareUnicodeCharacter{2228}{\ensuremath{\vee}}        % ∨
\DeclareUnicodeCharacter{2229}{\ensuremath{\cap}}        % ∩
\DeclareUnicodeCharacter{222A}{\ensuremath{\cup}}        % ∪
\DeclareUnicodeCharacter{2245}{\ensuremath{\cong}}       % ≅
\DeclareUnicodeCharacter{2248}{\ensuremath{\approx}}     % ≈
\DeclareUnicodeCharacter{2260}{\ensuremath{\neq}}        % ≠
\DeclareUnicodeCharacter{2264}{\ensuremath{\leq}}        % ≤
\DeclareUnicodeCharacter{2265}{\ensuremath{\geq}}        % ≥
\DeclareUnicodeCharacter{22C5}{\ensuremath{\cdot}}       % ⋅
\DeclareUnicodeCharacter{2295}{\ensuremath{\oplus}}      % ⊕
\DeclareUnicodeCharacter{2297}{\ensuremath{\otimes}}     % ⊗

% ---- Mathematics & symbols ------------------------------------------
\usepackage{amsmath,amssymb,bm,mathrsfs,mathtools}
\usepackage{amsthm}             % theorem environments (load BEFORE hyperref)

% ---- Theorem environments (shared counter set, TECT canonical) -------
% All theorems / lemmas / propositions / corollaries / remarks share a
% single counter so cross-references read consistently across a paper.
% Definitions get a separate counter (definitions are numbered in their
% own sequence by editorial convention).
\newtheorem{theorem}{Theorem}
\newtheorem{lemma}[theorem]{Lemma}
\newtheorem{proposition}[theorem]{Proposition}
\newtheorem{corollary}[theorem]{Corollary}

\theoremstyle{remark}
\newtheorem{remark}[theorem]{Remark}
\newtheorem{example}[theorem]{Example}

\theoremstyle{definition}
\newtheorem{definition}[theorem]{Definition}
\newtheorem{conjecture}[theorem]{Conjecture}
\newtheorem{hypothesis}[theorem]{Hypothesis}

% ---- Tables / floats / graphics --------------------------------------
\usepackage{booktabs}           % \toprule / \midrule / \bottomrule
\usepackage{array}
\usepackage{graphicx}
\usepackage{enumitem}           % \begin{itemize}[leftmargin=...] options
\usepackage{hyphenat}           % \hyp{} for breakable hyphens in \texttt
\usepackage{seqsplit}           % \seqsplit{} for long unbreakable strings
\sloppy                         % allow stretchier inter-word spacing to
                                % reduce overfull \hbox warnings on long URLs
                                % and long \texttt tokens (paragraph-level).
\emergencystretch=3em           % last-resort hbox stretch budget

% ---- Cross-references (hyperref last among the standard packages,
%      cleveref AFTER hyperref) ----------------------------------------
\usepackage[%
  colorlinks=true,%
  linkcolor=blue,%
  citecolor=blue,%
  urlcolor=blue,%
  breaklinks=true,%
  bookmarks=true,%
  bookmarksnumbered=true%
]{hyperref}
\usepackage[capitalise,nameinlink]{cleveref}

% ---- TECT-canonical author block macros -----------------------------
% (defined here so every paper has the same author / affiliation style)
\newcommand{\tectauthor}{%
  \author{Jusang Lee}%
  \email{jtkor@outlook.com}%
  \affiliation{Independent researcher, Republic of Korea}%
}

% ---- TECT-canonical archive footer ----------------------------------
\newcommand{\tectarchivefooter}{%
  \par\medskip
  \noindent\small\textit{Canonical archive}:
  \url{https://github.com/TECT-OpenPhysics/TECT}.\par%
}

% ---- TECT TODO / AUDIT-FLAG / HONEST-SCOPE inline marks --------------
% Used in drafts to highlight pending audit items without disturbing
% the body text style. Suppress via \tectsuppressmarks (define empty).
\providecommand{\tectauditflag}[1]{\textcolor{red}{\textbf{[AUDIT-FLAG: #1]}}}
\providecommand{\tecttodo}[1]{\textcolor{orange}{\textbf{[TODO: #1]}}}
\providecommand{\tecthonestscope}[1]{\textit{Honest scope: #1.}}

% =====================================================================
% End of TECT canonical preamble
% =====================================================================

%   \begin{document}
%   ...
%   \end{document}
%
% Build (every paper): `pdflatex Paper-NN && pdflatex Paper-NN`
% (the second pass resolves \ref{} / \cref{} forward references).
% A bash/PowerShell wrapper that automates this is at
% Docs/papers/papers/_shared/build-paper.sh / build-paper.ps1.
% =====================================================================

% ---- Document class --------------------------------------------------
% PRD class option of revtex4-2 with [reprint,onecolumn] gives the same
% APS heading / by-line / abstract style as PRL but in a wider single
% column that handles long display equations and theorem environments
% without column-break artefacts. nofootinbib keeps citations out of
% footnotes. The 11pt option restores standard font metrics; with T1
% fontenc + lmodern this gives non-bitmap glyphs.
\documentclass[%
  aps,prd,reprint,onecolumn,nofootinbib,11pt,%
  amsmath,amssymb,floatfix,longbibliography%
]{revtex4-2}

% ---- Encoding & fonts ------------------------------------------------
\usepackage[T1]{fontenc}        % proper 8-bit font encoding (no font corruption)
\usepackage[utf8]{inputenc}     % allow UTF-8 source
\usepackage{lmodern}            % scalable Latin Modern (replaces bitmap CM)
\usepackage{textcomp}           % extra text symbols
\usepackage{microtype}          % subtle kerning improvements

% ---- UTF-8 → LaTeX character mappings --------------------------------
% Maps the Unicode characters that occur in TECT body text (legacy from
% the chat-content auto-archival rule, CLAUDE.md §4) to LaTeX-safe
% commands. Without these declarations, T1+inputenc rejects the chars
% with "Unicode character X not set up for use with LaTeX".
% Adding declarations centrally (here) is preferable to per-file edits.
\DeclareUnicodeCharacter{00A7}{\S}              % §  (section sign)
\DeclareUnicodeCharacter{00B2}{\ensuremath{^{2}}}        % ²
\DeclareUnicodeCharacter{00B3}{\ensuremath{^{3}}}        % ³
\DeclareUnicodeCharacter{00B5}{\ensuremath{\mu}}         % µ (micro)
\DeclareUnicodeCharacter{00B7}{\ensuremath{\cdot}}       % · (middle dot)
\DeclareUnicodeCharacter{00D7}{\ensuremath{\times}}      % ×
\DeclareUnicodeCharacter{00E4}{\"a}             % ä
\DeclareUnicodeCharacter{00E9}{\'e}             % é
\DeclareUnicodeCharacter{00F6}{\"o}             % ö
\DeclareUnicodeCharacter{00FC}{\"u}             % ü
\DeclareUnicodeCharacter{010C}{\v{C}}           % Č
\DeclareUnicodeCharacter{010D}{\v{c}}           % č
\DeclareUnicodeCharacter{0394}{\ensuremath{\Delta}}      % Δ
\DeclareUnicodeCharacter{03A3}{\ensuremath{\Sigma}}      % Σ
\DeclareUnicodeCharacter{03A9}{\ensuremath{\Omega}}      % Ω
\DeclareUnicodeCharacter{03B1}{\ensuremath{\alpha}}      % α
\DeclareUnicodeCharacter{03B2}{\ensuremath{\beta}}       % β
\DeclareUnicodeCharacter{03B3}{\ensuremath{\gamma}}      % γ
\DeclareUnicodeCharacter{03B4}{\ensuremath{\delta}}      % δ
\DeclareUnicodeCharacter{03B5}{\ensuremath{\varepsilon}} % ε
\DeclareUnicodeCharacter{03BA}{\ensuremath{\kappa}}      % κ
\DeclareUnicodeCharacter{03BB}{\ensuremath{\lambda}}     % λ
\DeclareUnicodeCharacter{03BC}{\ensuremath{\mu}}         % μ
\DeclareUnicodeCharacter{03BD}{\ensuremath{\nu}}         % ν
\DeclareUnicodeCharacter{03C0}{\ensuremath{\pi}}         % π
\DeclareUnicodeCharacter{03C1}{\ensuremath{\rho}}        % ρ
\DeclareUnicodeCharacter{03C3}{\ensuremath{\sigma}}      % σ
\DeclareUnicodeCharacter{03C4}{\ensuremath{\tau}}        % τ
\DeclareUnicodeCharacter{03C6}{\ensuremath{\varphi}}     % φ
\DeclareUnicodeCharacter{03C7}{\ensuremath{\chi}}        % χ
\DeclareUnicodeCharacter{03C8}{\ensuremath{\psi}}        % ψ
\DeclareUnicodeCharacter{03C9}{\ensuremath{\omega}}      % ω
\DeclareUnicodeCharacter{2013}{--}              % – (en dash)
\DeclareUnicodeCharacter{2014}{---}             % — (em dash)
\DeclareUnicodeCharacter{2018}{`}               % ‘
\DeclareUnicodeCharacter{2019}{'}               % ’
\DeclareUnicodeCharacter{201C}{``}              % “
\DeclareUnicodeCharacter{201D}{''}              % ”
\DeclareUnicodeCharacter{2070}{\ensuremath{^{0}}}        % ⁰
\DeclareUnicodeCharacter{2071}{\ensuremath{^{i}}}        % ⁱ
\DeclareUnicodeCharacter{2122}{\textsuperscript{TM}}     % ™
\DeclareUnicodeCharacter{2126}{\ensuremath{\Omega}}      % Ω (ohm sign)
\DeclareUnicodeCharacter{210F}{\ensuremath{\hbar}}       % ℏ
\DeclareUnicodeCharacter{2115}{\ensuremath{\mathbb{N}}}  % ℕ
\DeclareUnicodeCharacter{2102}{\ensuremath{\mathbb{C}}}  % ℂ
\DeclareUnicodeCharacter{211D}{\ensuremath{\mathbb{R}}}  % ℝ
\DeclareUnicodeCharacter{2124}{\ensuremath{\mathbb{Z}}}  % ℤ
\DeclareUnicodeCharacter{211A}{\ensuremath{\mathbb{Q}}}  % ℚ
\DeclareUnicodeCharacter{2192}{\ensuremath{\to}}         % →
\DeclareUnicodeCharacter{21D2}{\ensuremath{\Rightarrow}} % ⇒
\DeclareUnicodeCharacter{2200}{\ensuremath{\forall}}     % ∀
\DeclareUnicodeCharacter{2203}{\ensuremath{\exists}}     % ∃
\DeclareUnicodeCharacter{2208}{\ensuremath{\in}}         % ∈
\DeclareUnicodeCharacter{2209}{\ensuremath{\notin}}      % ∉
\DeclareUnicodeCharacter{2227}{\ensuremath{\wedge}}      % ∧
\DeclareUnicodeCharacter{2228}{\ensuremath{\vee}}        % ∨
\DeclareUnicodeCharacter{2229}{\ensuremath{\cap}}        % ∩
\DeclareUnicodeCharacter{222A}{\ensuremath{\cup}}        % ∪
\DeclareUnicodeCharacter{2245}{\ensuremath{\cong}}       % ≅
\DeclareUnicodeCharacter{2248}{\ensuremath{\approx}}     % ≈
\DeclareUnicodeCharacter{2260}{\ensuremath{\neq}}        % ≠
\DeclareUnicodeCharacter{2264}{\ensuremath{\leq}}        % ≤
\DeclareUnicodeCharacter{2265}{\ensuremath{\geq}}        % ≥
\DeclareUnicodeCharacter{22C5}{\ensuremath{\cdot}}       % ⋅
\DeclareUnicodeCharacter{2295}{\ensuremath{\oplus}}      % ⊕
\DeclareUnicodeCharacter{2297}{\ensuremath{\otimes}}     % ⊗

% ---- Mathematics & symbols ------------------------------------------
\usepackage{amsmath,amssymb,bm,mathrsfs,mathtools}
\usepackage{amsthm}             % theorem environments (load BEFORE hyperref)

% ---- Theorem environments (shared counter set, TECT canonical) -------
% All theorems / lemmas / propositions / corollaries / remarks share a
% single counter so cross-references read consistently across a paper.
% Definitions get a separate counter (definitions are numbered in their
% own sequence by editorial convention).
\newtheorem{theorem}{Theorem}
\newtheorem{lemma}[theorem]{Lemma}
\newtheorem{proposition}[theorem]{Proposition}
\newtheorem{corollary}[theorem]{Corollary}

\theoremstyle{remark}
\newtheorem{remark}[theorem]{Remark}
\newtheorem{example}[theorem]{Example}

\theoremstyle{definition}
\newtheorem{definition}[theorem]{Definition}
\newtheorem{conjecture}[theorem]{Conjecture}
\newtheorem{hypothesis}[theorem]{Hypothesis}

% ---- Tables / floats / graphics --------------------------------------
\usepackage{booktabs}           % \toprule / \midrule / \bottomrule
\usepackage{array}
\usepackage{graphicx}
\usepackage{enumitem}           % \begin{itemize}[leftmargin=...] options
\usepackage{hyphenat}           % \hyp{} for breakable hyphens in \texttt
\usepackage{seqsplit}           % \seqsplit{} for long unbreakable strings
\sloppy                         % allow stretchier inter-word spacing to
                                % reduce overfull \hbox warnings on long URLs
                                % and long \texttt tokens (paragraph-level).
\emergencystretch=3em           % last-resort hbox stretch budget

% ---- Cross-references (hyperref last among the standard packages,
%      cleveref AFTER hyperref) ----------------------------------------
\usepackage[%
  colorlinks=true,%
  linkcolor=blue,%
  citecolor=blue,%
  urlcolor=blue,%
  breaklinks=true,%
  bookmarks=true,%
  bookmarksnumbered=true%
]{hyperref}
\usepackage[capitalise,nameinlink]{cleveref}

% ---- TECT-canonical author block macros -----------------------------
% (defined here so every paper has the same author / affiliation style)
\newcommand{\tectauthor}{%
  \author{Jusang Lee}%
  \email{jtkor@outlook.com}%
  \affiliation{Independent researcher, Republic of Korea}%
}

% ---- TECT-canonical archive footer ----------------------------------
\newcommand{\tectarchivefooter}{%
  \par\medskip
  \noindent\small\textit{Canonical archive}:
  \url{https://github.com/TECT-OpenPhysics/TECT}.\par%
}

% ---- TECT TODO / AUDIT-FLAG / HONEST-SCOPE inline marks --------------
% Used in drafts to highlight pending audit items without disturbing
% the body text style. Suppress via \tectsuppressmarks (define empty).
\providecommand{\tectauditflag}[1]{\textcolor{red}{\textbf{[AUDIT-FLAG: #1]}}}
\providecommand{\tecttodo}[1]{\textcolor{orange}{\textbf{[TODO: #1]}}}
\providecommand{\tecthonestscope}[1]{\textit{Honest scope: #1.}}

% =====================================================================
% End of TECT canonical preamble
% =====================================================================


% Paper-05-specific theorem-like environment (Principle) sharing the
% TECT shared theorem counter.
\newtheorem{principle}[theorem]{Principle}

\begin{document}

\title{Pillar~5 of the Topological Energy Condensate Theory:
chirality from shell-protected Dirac genericity, opposite-valley
pairing, mass protection, and the reduced Weyl/Dirac witness
theorem on the BCC condensate}

\author{Jusang Lee}
\email{jtkor@outlook.com}
\affiliation{Independent researcher, Republic of Korea}

\date{\today}

\begin{abstract}
We present the Pillar-5 chirality programme of the Topological
Energy Condensate Theory (TECT). The candidate mechanism
combines five layered results from the Math10--14 + Math171-AddA
canonical archive:
(i) Math10 shell-genericity --- once the Bragg mass is tuned away,
emergent Dirac propagation is generic in the symmetry-allowed
tensor space (Single-Channel and Two-Channel Genericity Theorems +
Dirac Genericity Principle);
(ii) Math11 minimal-model uniqueness --- five microscopic
constraints (finite-$q$ shell instability, quartic stabilisation,
internal $\mathbb{CP}^2$ geometry, enriched shell-Dirac mechanism,
second-order time dynamics) force the minimal field content to be
the complex triplet;
(iii) Math12 mass protection --- on the doubled block (after
opposite-valley pairing) the classification of constant Hermitian
operators identifies a four-dimensional symmetry-preserving space,
of which only two elements are genuine Dirac masses; the residual
valley charge generator forbids the remaining channels;
(iv) Math13 local Dirac block reduction --- under the residual
little-group mismatch, the local Dirac carrier reduces to
$\mathcal{H}_D \cong \mathbb{C}^2_{\rm valley} \otimes \mathbb{C}^2_{\rm int}$,
reducing the first remaining Dirac-closure condition to a finite
microscopic audit;
(v) Math14 reduced Weyl/Dirac witness theorem --- three theorems
(linearised quartic no-go, exact shell cancellation of the
isotropic first-order term at $q_*^2 = -Z/(2Y)$, locked cross-shell
quartic vanishing at zeroth order) close the reduced Weyl/Dirac
witness inside the Class II seed space;
(vi) Math171-AddA Hirzebruch--Riemann--Roch index formula (Top-impact
paper TI-1): $\mathrm{ind}(D_E^c) = r - \mu$ on $\mathbb{CP}^2$,
specialising to $\mathrm{ind} = 16$ at the rank-16 Standard-Model
fermion sector with $\mu = 0$ (NOT the previously incorrect
$\mathrm{Index}(D) = \mathrm{wind}(m)$ identification on the 3D
Brillouin torus, which would have required an unavailable
Callias / domain-wall framework).
The combined statement is a Pillar-5 \emph{candidate mechanism note}
at the canonical \texttt{Docs/status/TOE-FACT-SHEET.md} tier, NOT a
closed theorem of the form ``the chiral family count equals the
winding number''. A theorem-level rewrite within a Callias /
Niemi--Semenoff / domain-wall framework on the appropriate
non-compact extension remains an open programme item per
\texttt{Q-2026-05-02-Pillar5-Mechanism-Rewrite}.
\end{abstract}

\maketitle

% =====================================================================
\section{Introduction and the structural question}\label{sec:intro}
% =====================================================================

Chirality --- the asymmetric coupling of left- versus right-handed
fermions to weak-scale gauge interactions --- is one of the most
striking features of the Standard Model. The TECT framework
addresses the candidate \emph{origin} of chirality: rather than
postulating left-handed Weyl spinors as fundamental input fields, the
candidate proposal is that chirality emerges from the
topological-defect sector of the BCC condensate via shell-protected
Dirac couplings, opposite-valley pairing, and a residual-valley-symmetry
mass protection.

The present paper reports the Pillar-5 chirality programme as a
\emph{candidate mechanism note} (Rev~3, 2026-05-03) compiling the
canonical Math10--14 + Math171-AddA content. We do NOT claim a
closed theorem of the form ``$\mathrm{Index}(D) = \mathrm{wind}(m)$''
on the 3D Brillouin torus: such an identification would require a
Callias / Niemi--Semenoff / domain-wall framework that is not
provided by the present archive and is tracked separately as an open
follow-up
(\texttt{Q-2026-05-02-Pillar5-Mechanism-Rewrite}).

The structural question is therefore: \emph{what features of the
BCC condensate machinery are necessary, what is sufficient, and what
remains conjectural}? The paper organises the five Math10--14
results into a layered candidate mechanism (\S\S\ref{sec:Math10}--\ref{sec:Math14}),
relates the rank-16 family-count specialisation to the corrected
HRR formula of Math171-AddA / Top-impact paper TI-1
(\S\ref{sec:Math171-AddA}), and states the honest scope explicitly
in \S\ref{sec:scope}.

% =====================================================================
\section{Math10 shell-genericity: generic nonvanishing of
shell-protected Dirac couplings}\label{sec:Math10}
% =====================================================================

The first layer establishes that, once the Bragg mass at the locked
shell is tuned away, the emergence of Dirac-type linear dispersion in
the shell-protected band is generic in the symmetry-allowed
microscopic tensor space.

\begin{theorem}[Single-channel Dirac genericity, after Math10]
\label{thm:single-channel-genericity}
Let $\mathcal{T}_{\rm sym}$ be the linear space of microscopic
gradient-linear tensor couplings of the form
$N_{ij}^a$, where $i$ is the output spatial index, $j$ the gradient
index, and $a$ the internal channel label, restricted to the
$O_h$-equivariant subspace allowed by the Brazovskii operating-point
symmetries. Then the locus inside $\mathcal{T}_{\rm sym}$ on which
the shell-protected Dirac velocity in any single channel vanishes is
of strictly positive codimension; consequently the velocity is
\emph{generic} (i.e., non-vanishing on a Zariski-dense open subset).
\end{theorem}

\begin{theorem}[Simultaneous two-channel genericity, after Math10]
\label{thm:two-channel-genericity}
Under the same symmetry conditions, the simultaneous vanishing of
the Dirac velocities in two distinct shell-protected channels is of
strictly positive codimension in $\mathcal{T}_{\rm sym}$;
consequently the two-channel Dirac structure is generic.
\end{theorem}

\begin{principle}[Dirac Genericity Principle, after Math10]
\label{prn:Dirac-genericity}
Once the Bragg mass is tuned away and the microscopic enriched
tensors are not symmetry-forbidden, emergent Dirac propagation is
generic in the symmetry-allowed tensor space.
\end{principle}

\begin{remark}[Status of the genericity theorems]
The Math10 theorems are \emph{transversality / dimension-count}
results in the space of allowed microscopic couplings; they do NOT
by themselves construct the Standard-Model gauge structure or the
chiral family count. They establish that the Dirac mechanism is not
forbidden by any obvious symmetry obstruction.
\end{remark}

% =====================================================================
\section{Math11 minimal-model uniqueness: the complex
triplet}\label{sec:Math11}
% =====================================================================

The second layer addresses the question: \emph{which microscopic
field content is forced} by the conjunction of the Brazovskii
shell-instability mechanism, quartic stabilisation, internal
geometric structure, and the Math10 shell-Dirac mechanism?

\begin{theorem}[Math11 minimal-model uniqueness; informal statement]
\label{thm:Math11-minimal-model}
Among candidate microscopic models compatible with the conjunction
of:
\begin{enumerate}
\item Constraint~I --- finite-$q$ shell instability with locked
shell radius $q_0$ from Brazovskii kinetics;
\item Constraint~II --- quartic stabilisation with $g > 0$ ensuring
energy boundedness from below;
\item Constraint~III --- internal $\mathbb{CP}^2$ geometric
structure forced by the locked-triple algebra;
\item Constraint~IV --- the enriched shell-Dirac mechanism
(Math10) requiring an enriched internal doublet plus
gradient-linear microscopic tensor couplings;
\item Constraint~V --- second-order time dynamics consistent with
relativistic kinematics at the IR scale;
\end{enumerate}
the minimal field content is the complex triplet
$\Psi : \mathbb{R}^3 \to \mathbb{C}^3$. Lower-complexity
alternatives (real scalar, single complex scalar, complex doublet)
fail at least one of Constraints~I--V, and higher-complexity
alternatives (complex quartet, etc.) are non-minimal.
\end{theorem}

\begin{remark}[Reading guide]
The Math11 theorem is a \emph{constraint-counting} statement on the
microscopic field content, given the Brazovskii + shell-Dirac +
internal-$\mathbb{CP}^2$ inputs. The reader should view it as the
TECT-internal answer to the question ``why this particular
field content?'', not as an external no-go argument against general
alternative theories.
\end{remark}

% =====================================================================
\section{Math12 mass protection: classification of constant Dirac
masses on the doubled block}\label{sec:Math12}
% =====================================================================

After opposite-valley pairing on the locked shell, the shell-relevant
local Dirac block is doubled into a $\mathbb{C}^2_{\rm valley}
\otimes \mathbb{C}^2_{\rm int}$ structure. Math12 classifies the
constant Hermitian operators on this doubled block compatible with
the residual symmetry.

\begin{lemma}[Classification of constant Dirac masses, after Math12]
\label{lem:Math12-mass-classification}
Let $K$ be a constant Hermitian operator on the doubled block
$\mathbb{C}^2_{\rm valley} \otimes \mathbb{C}^2_{\rm int}$ that
preserves the little-group action at the band-crossing point.
By Schur's lemma applied to the irreducible little-group
representation $\rho$ on each valley, the space of
symmetry-preserving constant operators is four-dimensional:
\[
K = \alpha\, \tau_0 \otimes \mathbb{1} + \beta\, \tau_3 \otimes
\mathbb{1} + \gamma\, \tau_1 \otimes \mathbb{1} + \delta\, \tau_2
\otimes \mathbb{1},
\]
where $\tau_i$ are Pauli matrices on the valley pair $(+G_*,
-G_*)$. Of these four parameters, only two ($\gamma, \delta$ in the
above basis) are genuine Dirac-mass candidates that open a chiral gap
on the doubled block; the remaining two ($\alpha, \beta$) are
energy-shift / valley-energy-splitting terms, NOT genuine Dirac
masses.
\end{lemma}

\begin{theorem}[Mass protection from residual valley symmetry,
after Math12]
\label{thm:Math12-protection}
The residual valley charge generator $Q_V := \tau_3 \otimes
\mathbb{1}$ commutes with the shell Dirac kinetic operator and
forbids the addition of the genuine Dirac-mass channels at the
constant-operator level. Consequently, on the doubled block, the
chiral zero modes are protected against the addition of constant
mass terms by the residual valley symmetry; their gap can only be
opened by symmetry-breaking interactions of higher order
(non-constant, non-symmetric) that are not generic in the allowed
microscopic-coupling space.
\end{theorem}

% =====================================================================
\section{Math13 local Dirac block reduction
$\mathcal{H}_D \cong \mathbb{C}^2_{\rm valley}
\otimes \mathbb{C}^2_{\rm int}$}\label{sec:Math13}
% =====================================================================

The third layer reduces the first remaining Dirac-closure condition
to a finite microscopic audit on the doubled block.

\begin{proposition}[Math13 local Dirac block reduction]
\label{prop:Math13-block}
Under the little-group mismatch obtaining at the locked Brazovskii
operating point, the actual local Dirac carrier (the shell-relevant
local block $\mathcal{H}_D$) is NOT forced by an ordinary 2D
little-group irreducible representation. Instead, the local Dirac
carrier reduces to the tensor product
\[
\mathcal{H}_D \;\cong\;
\mathbb{C}^2_{\rm valley} \otimes \mathbb{C}^2_{\rm int},
\]
where:
\begin{itemize}
\item $\mathbb{C}^2_{\rm valley}$ is the opposite shell pair $(+G_*,
-G_*)$;
\item $\mathbb{C}^2_{\rm int}$ encodes the residual grading,
parity, or chirality at each valley.
\end{itemize}
The first remaining full Dirac-closure condition is reduced to a
finite microscopic audit on this four-dimensional carrier.
\end{proposition}

\begin{remark}[On the absence of an Atiyah--Singer-type closed-form
formula]
\label{rem:no-AS-closed-form}
The reduction to the finite carrier $\mathcal{H}_D$ is NOT an
Atiyah--Singer-type identity of the form ``$\mathrm{Index}(D) =
\mathrm{topological\ invariant}$'' on the 3D Brillouin torus. Such
an identity would require a Callias / Niemi--Semenoff / domain-wall
framework on a non-compact extension or a band-projector with
controlled IR boundary; the present BCC framework does not by itself
provide this. The Math13 proposition is therefore a
\emph{block-reduction} statement, not an index identity.
\end{remark}

% =====================================================================
\section{Math14 reduced Weyl/Dirac witness theorem: three
cancellation results}\label{sec:Math14}
% =====================================================================

The fourth layer provides three explicit theorems closing the
reduced Weyl/Dirac witness inside the Class~II seed space at the
Brazovskii-locked operating point.

\begin{theorem}[Linearised quartic no-go, after Math14 Theorem~1]
\label{thm:Math14-quartic-nogo}
At linear order in the fluctuation $\delta\Psi$ around the locked
Brazovskii background $\Psi_0(x) = A\, e^{i G \cdot x}\, z_0$ (with
$|G| = q_*$, $q_*^2 = -Z/(2Y)$, $z_0^\dagger z_0 = 1$), the quartic
self-coupling
$g(\Psi^\dagger \Psi)\Psi$
does NOT directly generate any enriched first-order operator of the
form $(\partial_i \Psi_0)(-i \partial_j) \otimes s^a$. Consequently
the quartic does not directly produce an enriched first-order tensor
$N_{ij}^a$ at linear order.
\end{theorem}

\begin{theorem}[Exact shell cancellation of the isotropic
first-order term, after Math14 Theorem~2]
\label{thm:Math14-shell-cancellation}
For the minimal isotropic kinetic operator
$Z\nabla^2 - Y\nabla^4$, the shell-projected first-order transport
term in the envelope $\chi$ vanishes exactly at the Brazovskii shell
$q_*^2 = -Z/(2Y)$.
\end{theorem}

\begin{theorem}[Locked cross-shell quartic vanishing at zeroth
order, after Math14 Theorem~3]
\label{thm:Math14-cross-shell}
The locked cross-shell quartic projection in the shell-restricted
envelope theory vanishes at zeroth order in the small-detuning
expansion. Consequently, the cross-shell quartic does NOT contribute
to the enriched first-order tensor structure required for the
shell-Dirac mechanism at zeroth order.
\end{theorem}

\begin{corollary}[Reduced Weyl/Dirac witness closed in Class~II seed
space, after Math14 Corollary]
\label{cor:Math14-reduced-witness}
The conjunction of Theorem~\ref{thm:Math14-quartic-nogo},
Theorem~\ref{thm:Math14-shell-cancellation}, and
Theorem~\ref{thm:Math14-cross-shell} closes the reduced Weyl/Dirac
witness theorem inside the explicit Class~II seed space at the
Brazovskii-locked operating point: the obstructions to the
shell-Dirac mechanism at the relevant order are eliminated, leaving
the genericity arguments of Math10 (\S\ref{sec:Math10}) unblocked.
\end{corollary}

% =====================================================================
\section{Math171-AddA Hirzebruch--Riemann--Roch index formula and
the rank-16 specialisation}\label{sec:Math171-AddA}
% =====================================================================

The chiral family count is addressed not by a winding-number
identification on the 3D Brillouin torus, but by the
Hirzebruch--Riemann--Roch (HRR) index formula on the canonical Pillar-4
realisation base, which reduces to a rank specialisation of the
general spin-c index.

\begin{theorem}[Math171-AddA HRR index formula, see Top-impact
paper TI-1 for full statement]
\label{thm:Math171-AddA}
For a holomorphic vector bundle $E$ of rank $r$ on $\mathbb{CP}^2$
with $c_1(E) = 0$ and $c_2(E) = \mu H^2$ for some integer $\mu \in
\mathbb{Z}$ (where $H$ is the hyperplane class), the spin-c Dirac
index with the canonical spin-c structure is
\[
\mathrm{ind}(D_E^c) \;=\; r - \mu \;\in\; \mathbb{Z}.
\]
Specialising to the rank-16 SO(10) chiral-fermion sector $r = 16$
at the canonical realisation $\mu = 0$ (Math179--198 archive
internal reading) gives $\mathrm{ind}(D_E^c) = 16$, consistent with
one Standard-Model family per chiral irrep at the level of the
historical $\mathbb{CP}^2$-route algebraic count.
\end{theorem}

\begin{remark}[Replacement of the $\mathrm{Index}(D) =
\mathrm{wind}(m)$ identification]
\label{rem:no-wind-identification}
The previously incorrect identification ``$\mathrm{Index}(D) =
\mathrm{wind}(m)$'' on the 3D Brillouin torus, applied to a real
scalar mass field $m(\mathbf{k})$ without a Callias-type / domain-wall
framework, does not hold in general. The correct chiral-count
machinery in the TECT framework is the conjunction of (a) the Math10
shell-genericity theorems (\S\ref{sec:Math10}), (b) the Math12 mass
protection (\S\ref{sec:Math12}), (c) the Math13 local Dirac block
reduction $\mathcal{H}_D \cong \mathbb{C}^2_{\rm valley} \otimes
\mathbb{C}^2_{\rm int}$ (\S\ref{sec:Math13}), and (d) the
HRR-formula family-count specialisation
$\mathrm{ind}(D_E^c) = r - \mu$ at $r = 16$, $\mu = 0$
(Theorem~\ref{thm:Math171-AddA}). The historical Pillar-4
$\mathbb{CP}^2$-route framing is preserved in the companion
top-impact paper TI-1; the canonical Pillar-4 realisation programme
operates on $\Sigma_0 = \mathbb{P}^1 \times \mathbb{P}^1$ per
Math305 (see Paper~04 for the full Pillar-4 atomic theorem).
\end{remark}

% =====================================================================
\section{Honest scope: what Pillar~5 establishes and what remains
open}\label{sec:scope}
% =====================================================================

\paragraph{What the present paper establishes.}
At the canonical mechanism-note tier per the canonical
\texttt{Docs/status/TOE-FACT-SHEET.md}, the present paper compiles a
five-layer candidate Pillar-5 chirality mechanism: (i) shell-Dirac
genericity (Math10 Theorems~\ref{thm:single-channel-genericity}--\ref{thm:two-channel-genericity}
and Principle~\ref{prn:Dirac-genericity}); (ii) minimal-model
uniqueness forcing the complex triplet (Math11
Theorem~\ref{thm:Math11-minimal-model}); (iii) mass protection from
residual valley symmetry (Math12
Lemma~\ref{lem:Math12-mass-classification} and
Theorem~\ref{thm:Math12-protection}); (iv) local Dirac block
reduction $\mathcal{H}_D \cong \mathbb{C}^2_{\rm valley} \otimes
\mathbb{C}^2_{\rm int}$ (Math13
Proposition~\ref{prop:Math13-block}); (v) reduced Weyl/Dirac witness
theorems closing the Class~II seed space (Math14
Theorems~\ref{thm:Math14-quartic-nogo}--\ref{thm:Math14-cross-shell}
and Corollary~\ref{cor:Math14-reduced-witness}). The HRR family-count
specialisation $\mathrm{ind}(D_E^c) = r - \mu$ at $r = 16$, $\mu = 0$
(Math171-AddA Theorem~\ref{thm:Math171-AddA}) provides the
rank-16 Standard-Model fermion identification at the historical
$\mathbb{CP}^2$-route algebraic level.

\paragraph{What the present paper does NOT establish.}
The paper does NOT establish (a) a closed theorem of the form
``the chiral family count equals a winding number'' on the 3D
Brillouin torus --- this would require a Callias / Niemi--Semenoff
/ domain-wall framework on a non-compact extension or an explicit
band-projector with controlled IR boundary, which the present
archive does not provide; (b) full unconditional Pillar-5
closure --- the canonical tier is a candidate mechanism with
explicitly stated layered hypotheses; (c) the family-count derivation
on the canonical $\Sigma_0 = \mathbb{P}^1 \times \mathbb{P}^1$
Pillar-4 base --- this requires the $\Sigma_0$ spin-c index calculation
(separate Todd class) which is not derived here; (d) the
$\mathrm{Cl}(3)$ Clifford-algebra closure at the Brillouin-zone-boundary
Dirac node --- this remains conditional on the Math233 spinor
pair-kernel programme.

\paragraph{Open follow-ups.}
The principal open programme item is the
\texttt{Q\hyp{}2026\hyp{}05\hyp{}02\hyp{}Pillar5\hyp{}Mechanism\hyp{}Rewrite}
ticket: a theorem-level rewrite within a Callias-type /
Niemi--Semenoff / domain-wall framework that promotes the present
mechanism note to an explicit chiral-count theorem with a topological
invariant on the appropriate non-compact extension. This is
mathematical-research-level work (estimated 2--4 weeks per the
\texttt{OPEN-QUESTIONS.md} entry) and is the structural mitigation
for the previously incorrect $\mathrm{Index}(D) = \mathrm{wind}(m)$
over-claim.

\paragraph{Implications for downstream pillars.}
The Pillar-5 mechanism feeds into Pillar~6 (generations and SM
embedding) at the rank-16 specialisation level via
Theorem~\ref{thm:Math171-AddA} (one SO(10) $\mathbf{16}$ per chiral
irrep). The mass protection from valley symmetry
(Theorem~\ref{thm:Math12-protection}) supports the candidate
Standard-Model neutrino-mass programme at the leading-order level;
explicit derivation of neutrino masses from the
$\mathbb{C}^2_{\rm int}$ valley structure is an open programme
item. The Pillar-7 quantum-consistency programme (Ward identities,
SO(10) anomaly cancellation) is supported by the family-count
specialisation but does not in itself force the family count to
exactly three; the rank-16 $\to$ three-generation question remains
the subject of the companion Paper~07 (which is itself currently
\texttt{[EXTERNAL\_USE\_FORBIDDEN]} pending the
Q-2026-05-02-Pillar7-Three-Generation-Rewrite follow-up).

% =====================================================================
\begin{acknowledgments}
The author thanks the Anthropic Claude Opus 4.7 collaborator for
the synthesis of the Pillar-5 canonical content from the Math10--14
shell-genericity + minimal-model + mass-protection + Dirac-block-reduction
+ reduced-Weyl/Dirac-witness chain plus Math171-AddA HRR formula
into the present revision (Rev~3, 2026-05-03), and for the Math314
family hostile-referee audit cycle that drove the prior Rev~1/Rev~2
corrections (in particular the removal of the incorrect
$\mathrm{Index}(D) = \mathrm{wind}(m)$ identification on the 3D
Brillouin torus). This work is part of TECT, canonical archive at
\url{https://github.com/TECT-OpenPhysics/TECT}.
\end{acknowledgments}

% =====================================================================
\bibliographystyle{apsrev4-2}
\bibliography{../_shared/tect-references}
% =====================================================================

\end{document}
